\newpage
\section{Theoretischer Hintergrund}\label{theoretischer_hintergrund}
% Einleitender Absatz in das Kapitel.


\subsection{Ökologischer Hintergrund}\label{okologischer_hintergrund}

Als Heiden werden baumfreie Offenlandschaften auf nährsoffarmen, sauren Böden bezeichnet, die überwiegend vppn niedrigwüchsigen Zwergsträuchern aus der Familie der Heidekrautgewächse (Ericaceae) bewachsen sind. Prägende und namensgebende Leitart ist dabei die Besenheide (\textit{Calluna vulgaris}) \cite{Webb.1998}.

Die heutigen Heideflächen haben sich vr etwa 4.000 Jahren großflächig als Folge menschlicher Waldrodungen entwickelt \cite{Webb.1998}. Die brachliegenden Flächen wurden anschließend auf vielfältige Weise genutzt, etwa zur Viehhaltung, zum Torfabbau, für die Brandrodung, zur Gewinnung von Brennstoff sowie als Fläche für Tierfutter. Vor allem die Beweidung hielt die Böden nährstoffarm und verhinderte zugleich, dass sich auf den gerodeten Flächen erneut Wald entwickeln konnte, sodass sich dadurch die Heide ausbreiten konnte. Heidekraut selbst gab es allerdings bereits vor dem menschlichen Einfluss, wenn auch nur als Beimischung auf kargen Waldböden, dort wo ein lichtes, offenes Kronendach genügend Licht auf den Boden fallen ließ. Damit zählen Heideflächen zu den wichtigsten Kulturlandschaften Westeuropas, also zu den durch menschliche Nutzung über lange Zeit geprägten und gestalteten Landschaften - im Gegensatz zur vom Menschen unberührten Naturlandschaft. Erstreckten sie sich einst über mehrere Millionen Hektar, sind heute nur noch etwa 350.000 Hektar übrig. Gerade wegen dieser starken Verluste gelten Heidegebiete als Flächen mit hohem Naturschutzpotential \cite{Webb.1998}.

Bleibt die zuvor beschriebene Nutzung aus, setzt an den Heidestandorten eine Sukzession ein, also die gerichtete zeitliche Abfolge von Pflanzengesellschaften an einem Standort - hier von Zwergstrauchheide über eine zunehmende Vergrasung bis hin zu Gehölzen und langfristig zurück zum Wald. Ohne Pflege wachsen nach und nach Pioniergehölze auf, etwas Schlehen, Rosen- und Weißdornarten sowie Bäume wie Birken, Zitterpappeln und Kiefern \cite{Schoknecht.1998}. Diese Dominanzbestände überwuchern die Heide, indem sie ihr Nährstoffe entziehen und beschatten, sodass die niedrigwüchsigen Zwergsträucher zu wenig Licht erhalten. Neben der ausbleibenden Pflege ist die atmosphärische Stickstoffdeposition der zentrale Treiber der heutigen Degradation. Über Jahrzehnte gelangt Stickstoff aus Landwirtschaft und Verkehr über die Luft in die Heide und reichert die von Natur aus nährstoffarmen Böden schleichend an, wobei die Einträge vielerorts über der ökologischen Belastungsgrenze liegen \cite{Bobbink.2010}. Dadurch verschiebt sich die Konkurrenz zugunsten schnellwüchsiger Gräser:  Auf Trockenheiden wird die Besenheide (\textit{Calluna vulgaris}) durch das Drahtschmielengras (\textit{Deschampsia flexuosa}) verdrängt, auf Feuchtheiden \textit{Erica tetralix} durch das Pfeifengras (\textit{Molinia caerulea}). Grund dafür ist, dass die Gräser den zusätzlichen Stickstoff effizienter aufnehmen und binden. Mit steigender Stickstoffverfügbarkeit wird die Konkurrenz zunehmend asymmetrisch, und die Zwergsträucher werden unterdrückt \cite{Friedrich.2011}. 

Beide Prozesse - die schleichende Vergrasung wie die anschließende Ausbreitung von Gehölzen - laufen letztlich in der Verbuschung zusammen. Unter Verbuschung versteht man die zunehmende Ausbreitung von Sträuchern und Gehölze auf der eigentlich offenen Heidefläche. Der Verbuschungsgrad bezeichnet das quantitative Maß dieser Ausbreitung, also den Anteil der von Gehölzen bedeckten Fläche. Da die Verbuschung die charakteristische Zwergstrauchvegetation verdrängt und die offene Struktur der Heide auflöst, gilt sie als Degradationsindikator für den Erhaltungszustand von Heideflächen \cite{Schmidt.2017}. Der Erhaltungszustand wiederum ist die Bewertungsgröße, anhand derer die Flora-Fauna-Habitat-Richtlinie den Zustand der geschützten Lebensräume erfasst \cite{RatEG.1992}. Ein hoher Verbuschungsgrad steht damit für einen ungünstigen Erhaltungszustand der Heide.

Um die Heide offenzuhalten und der Sukzession entgegenzuwirken, war schon früher eine kontinuierliche Nutzung nötig, aus der die heutigen Pflegemaßnahmen hervorgegangen sind. So wurde durch das Plaggen die nährstoffreiche Humusschicht samt Heidevegetation vom Oberboden abgetragen. Die Plaggen dienten als Einstreu in den Ställen und gelangten später als Dünger auf die Äcker \cite{Niemeyer.2007}. Mit der Einführung von Mineraldünger und ausreichend verfügbarem Stroh wurde diese Methode jedoch überflüssig. Auch die weit verbreitete Schafzucht hielt Bäume und Sträucher zurück, da die Schafe die jungen Triebe abfressen \cite{Bakker.1983}. Sie ging zurück, als Baumwolle und billige Wollimporte aus Übersee aufkamen. Daneben wurde die Mahd eingesetzt, bei der die oberen Pflanzenschichten mit Sense oder Sichel abgeschnitten und der Bestand so verjüngt wurde, wobei zugleich junge Baumtriebe und Sträucher entfernt wurden. Da diese traditionellen Nutzungsformen weggefallen sind, muss der Naturschutz die Offenhaltung heute mit gezielten Pflegemaßnahmen übernehmen. Der maschinelle Plaggenhieb bzw. das Schoppern entfernt Boden- und Rohhumusschichten und entzieht dem System damit gezielt Stickstoff. Die abgetragenen Schichten müssen vollständig aus dem Gebiet entfernt werden, da sich sonst konkurrenzstarke Gräser wie das Land-Reitgras wieder ausbreiten \cite{Niemeyer.2007}. Bei der Beweidung mit Schafe oder Heidschnucken ist das nächtliche Aufstallen entscheidend, also das Zusammenpferchen auf Äckern oder im Stall, damit der Dung nicht in die Heide zurückgelangt, sondern direkt auf dem Acker landet und die Fläche nährstoffarm bleibt \cite{Bakker.1983,RosaGarcia.2013}. Die Mahd zählt weiteren zu den gängigsten Methoden, wobei der Abtransport des Mähguts unbedingt notwendig ist, um eine erneute Nährstoffanreicherung zu vermeiden \cite{Haerdtle.2009}. Das Entkusseln geht noch einen Schritt weiter und entfernt aufkommende Gehölze wie Kiefer oder Birke rein mechanisch durch Absägen und Zurückschneiden. Schließlich wird auch das kontrollierte Abbrennen kleiner Flächen angewendet, das dem Rohhumus Nährstoffe entzieht und offene, konkurrenzarme Standorte für das Heidekraut schafft. Anders als beim Plaggen bleibt der Phosphor dabei jedoch über die Asche im System \cite{Mohamed.2007}. Vergleichende Untersuchungen zeigen, dass alle diese Verfahren dem System Nährstoffe entziehen, bei anhaltend hoher Stickstoffdeposition aber in kurzen Abständen von etwa fünf Jahren wiederholt werden müssen \cite{Haerdtle.2009}.

Diese Pflegemaßnahmen sind personal-, kosten- und zeitintensiv und müssen regelmäßig wiederholt werden. Ihr gezielter Einsatz und die Kotrolle ihres Erfolgs setzen daher voraus, dass sich der Verbuschungsgrad großflächig, objektiv und wiederholbar bestimmen lässt - Anforderungen, die eine rein feldbasierte Kartierung nur schwer erfüllt \cite{Kissling.2024}. Damit rücken fernerkundliche und automatisierte Auswertungsverfahren in den Fokus, deren methodische Grundlagen die folgenden Abschnitte behandeln: von der Rekonstruktion dreidimensionaler Punktwolken mittels Structure from Motion bis zu ihrer semantischen Segmentierung mit tiefen neuronalen Netzen.

\subsection{Structure from Motion}\label{structure_from_motion}

\subsubsection{Allgemeines}\label{sfm_allgemeines}
% Grundlagen zur Struktur für Bewegung. Beispielzitat: \footcite[\vglf][\pagef 1]{Balzert.2008}

- Ausgangslage für semantische Segmentierung sind aufgenommene Drohnenbilder im RGB Spektrum (2D)
-für meinen Teil der Arbeit wurden die Bilder in SfM Punktwolken mittels Software (welche gneua noch nachfragen) umgerechnet

- wohingegen LiDAR Punktwolken direkt mit einem extra dafür vorgesehenen Sensor am UAV aufgenommen werden und dann nur noch mittles einer firmeneigenen Software in das passende Format (z.B. LAS) konvertiert werden muss
- Structure from Motion entsthet dabei durch spezielle Algortihmen asu eigentlich 2D RGB Bildern
-SfM Algorithmen kombinieren viele stark überlappende Bildframes zu 2D Bildmosaiken und zu 3D Punktwolken und nutzen in UAV-Bildern entahltene photogrammeteische Informationen um hochaufgelösste DTMs und DSMs zu erstellen. 
- für Erstellung von DTMs zwar viele verschiedene Algorithmen zur Verfügung, aber Wahl der Parameter, damit verbundene Unsicherheiten oft nicht gut dokumentiert
- allerdings Open Source Softwaretools zur Verfügung 
- z.B. ColMap für einen reproduzierbaren Open-Source-Workflow für SfM Pipelones -> enthalten feature extraction -> feature matching -> geometric verification -> structure and motion reconstruction

allgemeine Beschreibung SfM Abgrenzung zu klassicher Photogrammtrie:

- SfM basiert auf denselben Grundprinzipien woie die stereoskopische Photogrmmatrie:
- 3D-Strukutr lässt sich aus einer Reihe überalppender, versetzter Bilder ableiten
- SfM Methode ermittelt Kamerapoition und  Szenegeometrie automatisch und gleichzeitig mithilfe eiens stark redundaten “bundle adjustment” auf Basis zugeordnenter Bildmerkmale

- bei klassicher photgrammetrie müssen vorab die Kameraposition und Kameraorinetoerung bekannt sein oder es müssen Paspunkte mit baknnten 3D Koordinaten exsistieren
- SfM braucht das nicht: Position und 3D Strukur allein aus Bildern automaitsch bestimmt

- stammt aus Computer Visiion Forschung in den 1990er jahren
- wurde bekannt durch die Arbeit von Snavely (2008) (beschreibt Ansätze von SfM im Cloud-Verarbeitungskontext von Microsoft Photosynth)

- erster Schritt: markante Bildpunkte in jedem Einzelbild finden z.B. mit SIFT
- diese “keypoints” sind invariant gegenüber Skalierung/rotation und osgar teilweise gegenüber Belcuhtungsänderungen
- zu jedem Keypoint wird ein Deskriptor berechnet, der ein einduetiges Zuordnen über viele Bilder hinweg erlaubt
-Qulaität/Dichte der späteren Puntwolke hängt von Bildtextur und UAflösung ab

- korrespondierende Merkmale müssen in mindesstens drei Bildern sichtbar sein
-hoher Überlappungsgrad und möglichst viele Bilder aus verschiedenen Postiione verbessern das Ergebis

-SfM Punktwolke zunächst in relativem „image space“-CRS ohne echten MAßtab
-Georeferenzierung erfolgt anschließend über wenig Passpunkte (GCPs) die über ein „3D similarity transform“ in reale „object-space Koordinaten überführt werden

-SfM liefert eine sehr dichte Punktwolke (bis <1000 Punkte/m²)
-anders als in geomorphologischen SfM-Anwendungen, die auf ein reines Geländemodell filtern, sind hier gerade die Nicht-Boden-Punkte (Vegetation, Gehölze) das Ziel der Segmentierung
% --- Muster für eine Abbildung (auskommentiert) ---
% \begin{figure}[H]
% \caption{Funktionsweise von Tesseract OCR}\label{fig:tesseract}
% \includegraphics[width=0.9\textwidth]{tesseract}
% \\
% Quelle: Eigene Darstellung
% \end{figure}

\subsection{Convolutional Neural Networks (CNNs)}\label{cnns}

dabei unterteilt er diese Mehtoden in vier Hauptkategorien, wobei die erste und zweite Kategorie eher großflächige, ökologische Zusammenhänge aus Satelliten- und Luftbildern analysieren, dringen die dritte und vierte in die Detaillebene, nutzen mehr auch LIDAR Daten und stellen die anatomischen Details einzelner Bäume dar 
- in der 1. KAT. wird eine "forest ressource assessment" (FRA) durchgeführt, wobei das CNN U-Net mittels semantischer Segmentierung zur Kartierung von Waldtypen am Beispiel des atlantischen Regenwaldes, eingesetzt wurde
- als 2. KAT wurde die "forest ecological disturbance" (FED) überwacht um mithilfe von YOLO-Modellen Waldbrandrauch frühzeitig zu erknnen und mithilfe von Mask R-CNN, kranke von einem Pilz befallene Bäume zu identifizieren
- die "forest parameter retrieval at the individual scale" (FPI) wurde angewendet und mittels Deeplab aus Drohnenbildern einzelne Buamkronen zu identifizieren
- außerdem kann man mit dieser Methode, über das CNN ResNet auch einzelne Baumarten genau identifizieren
- die vierte Hauptkategorie in der man Deep-Learning Anwendungen in Bezug auf Waldmonitoring unterteilt, ist die "tree organ-level measurement"(TOM)
- diese wird meist als Vorverarbeitungsschritt eingesetzt um dann noch detaillierter zu Untersuchen, da diese Methode besipielsweise mithilfe von LIDAR Daten die Holzstrukur und die Blätter eines Baumes trennen kann, was dann wiederum die Grundlage für die Bestimmung von Baumarten ist
- auch wird TOM vor allem im landwirtschaftlichen Bereich genutzt um Früchte zu lokalisieren und somit den Ernteertrag abzuschätzen

-man sieht klar wie viel mehr flexibel deep learning algorithmen im verlgeich zu klassichen machine learning algorithmen sind
- wohingegen diese noch stark auf die manuelle Merkmalsextraktion angewiesen sind und bei unregelmäßigen oder komplexen Datenmengen (insebsondere bei Baumstrukturen) schnell an ihre Grenzen stoßen
- im Vergleich dazu kann deep learning automatisch aus hochdimensionalen, nichtlinearen Rohdaten wichtige und entscheidene Merkmale extrahieren, lernen damit umzugehen und dann auch zu modellieren
- diese Flexibilität von deep learning zeigt sich insbesondere dadruch, dass die 4 verschiedenen Methoden die CNNs ganz flexibel nutzen und exakt an das Problem anpassen können
- die Pixel- und patch basierten CNNs wie U-Net oder ResNet kartieren fließende Übergänge von Waldtypen, wohingegen Objektbasierte CNNs wie YOLO oder Faster R-CNN, gezielt einzelne Organe wie Früchte, kranke Blätter oder Äste in komplexen Umgebungen lokalisieren


\subsubsection{Allgemeines}\label{cnns_Allgemeines}
% Grundlagen künstlicher neuronaler Netze.

\subsection{Geometric Deep Learning}\label{geometric_deep_learning}

\subsubsection{Semantische Segmentierung von 3D-Punktwolken}

\subsubsection{Beispiel: Pointnet}\label{pointnet}
% konkretes Beispiel für Geometrische Deep Learning (Vorgänger).

\subsubsection{Beispiel: Pointnet++}\label{pointnet++}
% konkretes Beispiel für Geometrische Deep Learning (aktueller Stand).

% \subsubsection{MLPs (Multilayer Perceptrons)}\label{mlps}



